El diseño experimental se estructuró a partir de un laboratorio físico controlado, concebido bajo una red de área local (LAN) completamente aislada del entorno institucional, lo cual permitió ejecutar simulaciones de ataques en un espacio confinado que eliminó interferencias externas y redujo de manera significativa la aparición de falsos positivos, al mismo tiempo que se resguardó la integridad de la red universitaria frente a posibles riesgos de seguridad derivados de la actividad maliciosa simulada.\\
En este contexto, la arquitectura de red adoptó una topología física en estrella en la que un servidor central de monitoreo y cuatro estaciones de ataque se interconectaron mediante un switch Ethernet y cableado UTP categoría 5e/6, configuración que aseguró una comunicación estable y de baja latencia y, a su vez, facilitó la captura íntegra del tráfico generado durante las interacciones ofensivas, siendo el servidor central una computadora de escritorio dedicada exclusivamente a alojar la plataforma de ciberseguridad T-Pot sobre una instalación limpia de Ubuntu Server 24.04.3 LTS, sistema operativo seleccionado por su estabilidad y compatibilidad nativa con tecnologías de contenedorización como Docker, sobre el cual se desplegó el conjunto de honeypots y herramientas de análisis integradas en dicha plataforma.\\
Por su parte, las estaciones de ataque correspondieron a cuatro computadoras portátiles independientes, configuradas con sistemas operativos heterogéneos (Linux y Windows) con el propósito de emular distintos perfiles de atacantes, de modo que cada estación fue preparada con herramientas y scripts especializados que permitieron ejecutar diversos vectores de ataque, entre los que se incluyeron escaneos de puertos mediante utilidades como Nmap, intentos de fuerza bruta contra servicios expuestos, particularmente SSH, así como pruebas de denegación de servicio (DoS) basadas en técnicas de inundación (flood).\\
Cabe señalar que, con el fin de preservar la seguridad perimetral de la red institucional, se implementó una red interna cerrada sin conexión a internet ni a otros segmentos universitarios, por lo que toda la actividad maliciosa registrada se originó exclusivamente desde direcciones IP locales, simulando un escenario de amenaza interna que permitió observar con mayor detalle patrones de comportamiento característicos, tales como la progresión desde fases de reconocimiento inicial mediante escaneo hasta etapas de explotación activa como la fuerza bruta o la ejecución de comandos, aislamiento que no solo cumplió con las normativas éticas y de seguridad establecidas por la institución, sino que además posibilitó una observación pasiva y prolongada de las tácticas empleadas por los atacantes, facilitando la recolección de evidencias forenses sin comprometer activos reales.\\
Finalmente, la configuración unificada de T-Pot, en conjunto con herramientas complementarias como Kibana y Htop, permitió correlacionar en tiempo real los eventos capturados, transformando los registros de tráfico y ataques en visualizaciones accesibles que facilitaron la identificación de tendencias, la intensidad de los ataques y el consumo de recursos durante las simulaciones, de manera que el escenario experimental no solo funcionó como un banco de pruebas para validar la eficacia de la plataforma en la detección de intrusiones, sino que también proporcionó un marco controlado y reproducible para el análisis de técnicas ofensivas contemporáneas, contribuyendo al fortalecimiento de estrategias proactivas de seguridad en entornos académicos y organizacionales.